\section*{Introduction}
\addcontentsline{toc}{section}{Introduction}

% TODO: Rédiger l'introduction du chapitre 2
% Points à couvrir:
% 1. Objectif du chapitre (Analyse et spécification des besoins)
% 2. Démarche suivie
% 3. Plan du chapitre

[CONTENU À REMPLIR: Introduction du chapitre 2 - Analyse et spécification des besoins]

\section{Identification des acteurs et des exigences}

% TODO: Identifier les acteurs du système
% Points à inclure:
% - Utilisateurs finaux
% - Administrateurs
% - Autres parties prenantes
% - Rôles et responsabilités

\subsection{Acteurs du système}

[CONTENU À REMPLIR: Identification des acteurs]

% Exemple:
% \begin{itemize}
%     \item Acteur 1: [description]
%     \item Acteur 2: [description]
%     \item Acteur 3: [description]
% \end{itemize}

\subsection{Exigences fonctionnelles}

% TODO: Lister les exigences fonctionnelles
% Points à inclure:
% - Fonctionnalités principales du système
% - Comportement attendu
% - Interactions avec les acteurs

[CONTENU À REMPLIR: Exigences fonctionnelles]

% Exemple:
% \begin{enumerate}
%     \item EF1: [description]
%     \item EF2: [description]
%     \item EF3: [description]
% \end{enumerate}

\subsection{Exigences non fonctionnelles}

% TODO: Lister les exigences non fonctionnelles
% Points à inclure:
% - Performance
% - Sécurité
% - Scalabilité
% - Disponibilité
% - Accessibilité
% - Maintenabilité

[CONTENU À REMPLIR: Exigences non fonctionnelles]

% Exemple:
% \begin{itemize}
%     \item ENF1 (Performance): [description]
%     \item ENF2 (Sécurité): [description]
%     \item ENF3 (Scalabilité): [description]
% \end{itemize}

\section{Modélisation des cas d'utilisation}

% TODO: Créer les diagrammes de cas d'utilisation
% Points à inclure:
% - Interactions entre acteurs et système
% - Scénarios principaux
% - Scénarios alternatifs

\subsection{Diagramme de cas d'utilisation global}

% TODO: Présenter le diagramme global
% Points à inclure:
% - Vue d'ensemble des cas d'utilisation
% - Acteurs impliqués
% - Relations entre cas d'utilisation

[CONTENU À REMPLIR: Description du diagramme de cas d'utilisation global]

% Exemple:
% \begin{figure}[H]
%     \centering
%     \includegraphics[width=0.8\textwidth]{images/use_case_diagram.png}
%     \caption{Diagramme de cas d'utilisation global}
%     \label{fig:use_case_global}
% \end{figure}

\subsection{Description détaillée des cas d'utilisation}

% TODO: Décrire chaque cas d'utilisation
% Points à inclure:
% - Nom du cas d'utilisation
% - Acteur principal
% - Condition préalable
% - Flux principal
% - Flux alternatifs
% - Condition postcondition

[CONTENU À REMPLIR: Description détaillée des cas d'utilisation]

% Exemple de cas d'utilisation:
% 
% \subsubsection{Cas d'utilisation: [Nom]}
% \begin{itemize}
%     \item \textbf{Acteur principal:} [Nom de l'acteur]
%     \item \textbf{Acteurs secondaires:} [Autres acteurs]
%     \item \textbf{Condition préalable:} [Conditions]
%     \item \textbf{Flux principal:} 
%     \begin{enumerate}
%         \item Étape 1
%         \item Étape 2
%         \item Étape 3
%     \end{enumerate}
%     \item \textbf{Flux alternatifs:} [Description]
%     \item \textbf{Condition postcondition:} [Conditions]
% \end{itemize}

\section{Modélisation des données}

% TODO: Concevoir le modèle de données
% Points à inclure:
% - Entités du système
% - Attributs
% - Relations entre entités
% - Cardinalités

\subsection{Diagramme de classes}

% TODO: Présenter le diagramme de classes UML
% Points à inclure:
% - Classes principales
% - Attributs et méthodes
% - Relations (association, héritage, etc.)

[CONTENU À REMPLIR: Description du diagramme de classes]

% Exemple:
% \begin{figure}[H]
%     \centering
%     \includegraphics[width=0.9\textwidth]{images/class_diagram.png}
%     \caption{Diagramme de classes}
%     \label{fig:class_diagram}
% \end{figure}

\subsection{Modèle entité-association (MCD/MLD)}

% TODO: Présenter le modèle conceptuel et logique des données
% Points à inclure:
% - Entités principales
% - Attributs
% - Relations
% - Clés primaires et étrangères

[CONTENU À REMPLIR: Modèle conceptuel et logique des données]

% Exemple:
% \begin{figure}[H]
%     \centering
%     \includegraphics[width=0.9\textwidth]{images/data_model.png}
%     \caption{Modèle entité-association}
%     \label{fig:data_model}
% \end{figure}

\section{Spécifications détaillées}

% TODO: Documenter les spécifications détaillées
% Points à inclure:
% - Flots de travail (workflows)
% - Règles métier
% - Contraintes
% - Interfaces utilisateur (wireframes/mockups)

\subsection{Workflows et processus métier}

[CONTENU À REMPLIR: Description des workflows et processus métier]

% Exemple:
% \begin{figure}[H]
%     \centering
%     \includegraphics[width=0.8\textwidth]{images/workflow.png}
%     \caption{Workflow principal}
%     \label{fig:workflow}
% \end{figure}

\subsection{Règles métier}

[CONTENU À REMPLIR: Énumération des règles métier]

% Exemple:
% \begin{enumerate}
%     \item Règle 1: [Description]
%     \item Règle 2: [Description]
%     \item Règle 3: [Description]
% \end{enumerate}

\section*{Conclusion}
\addcontentsline{toc}{section}{Conclusion}

% TODO: Conclure le chapitre 2
% Points à inclure:
% - Synthèse des besoins identifiés
% - Résumé du modèle de données
% - Transition vers la conception

[CONTENU À REMPLIR: Conclusion du chapitre 2]
